\documentclass{article}

\usepackage{amsmath,amssymb}
\usepackage{xurl}
\usepackage[hidelinks]{hyperref}
\setlength{\emergencystretch}{2em}

% Shared commands for notes in docs/paper-gaps/.

\DeclareUnicodeCharacter{2080}{\ifmmode{}_{0}\else\textsubscript{0}\fi}
\DeclareUnicodeCharacter{2081}{\ifmmode{}_{1}\else\textsubscript{1}\fi}
\DeclareUnicodeCharacter{2082}{\ifmmode{}_{2}\else\textsubscript{2}\fi}
\DeclareUnicodeCharacter{2099}{\ifmmode{}_{n}\else\textsubscript{n}\fi}

\newcommand{\Lean}{\textsc{Lean}}
\ExplSyntaxOn
\NewDocumentCommand{\leanid}{m}{
  \ifmmode
    \textsf{\detokenize{#1}}
  \else
    \group_begin:
    \urlstyle{sf}
    \tl_set:Nn \l_tmpa_tl {#1}
    \tl_replace_all:Nnn \l_tmpa_tl {\_} {_}
    \exp_args:NV \path \l_tmpa_tl
    \group_end:
  \fi
}
\ExplSyntaxOff
\providecommand{\tr}{\operatorname{tr}}
\newcommand{\tnleanIssueBase}{https://github.com/LionSR/TNLean/issues/}
\newcommand{\tnleanPrBase}{https://github.com/LionSR/TNLean/pull/}
\newcommand{\tnleanDocBase}{https://sirui-lu.com/TNLean/docs/}
\newcommand{\ghissue}[1]{\href{\tnleanIssueBase#1}{GitHub issue \##1}}
\newcommand{\ghpr}[1]{\href{\tnleanPrBase#1}{PR \##1}}
\newcommand{\doclink}[1]{\href{\tnleanDocBase#1.html}{\path{#1}}}

% Dirac notation and the identity operator, as in blueprint/src/macros/common.tex.
% \providecommand keeps note-local definitions authoritative.
\usepackage{dsfont}
\providecommand{\ket}[1]{|#1\rangle}
\providecommand{\bra}[1]{\langle #1|}
\providecommand{\braket}[2]{\langle #1 | #2 \rangle}
\providecommand{\ketbra}[2]{|#1\rangle\!\langle #2|}
\providecommand{\Id}{\mathds{1}}
\providecommand{\one}{\mathds{1}}
\providecommand{\C}{\mathbb{C}}
\providecommand{\MN}[1]{M_{#1}(\C)}
\providecommand{\MD}{M_D(\C)}

% External referencing for paper sources.  The build helper writes sanitized
% auxiliary files under build/paper-aux/ and excludes duplicated source labels.
\usepackage{xr}

% Import a generated paper auxiliary file with a caller-supplied label prefix.
% The candidate paths support builds from docs/paper-gaps/, the repository root,
% and build/paper-aux/ itself.
\newcommand{\importpaperaux}[2]{%
  \IfFileExists{../../build/paper-aux/#2.aux}{%
    \externaldocument[#1]{../../build/paper-aux/#2}%
  }{%
    \IfFileExists{build/paper-aux/#2.aux}{%
      \externaldocument[#1]{build/paper-aux/#2}%
    }{%
      \IfFileExists{#2.aux}{%
        \externaldocument[#1]{#2}%
      }{}%
    }%
  }%
}

% General external reference macros
\newcommand{\paperref}[2]{\ref{#1-#2}}
\newcommand{\papereqref}[2]{(\ref{#1-#2})}

% CPSV16 source (1606.00608 / MPDO-22-12-17-2.tex) external document and shortcuts
\importpaperaux{cpsv16-}{1606.00608/MPDO-22-12-17-2-tnlean}
\newcommand{\cpsvref}[1]{\paperref{cpsv16}{#1}}
\newcommand{\cpsveqref}[1]{\papereqref{cpsv16}{#1}}

% Verdict marker: \gapnote{<kind>}{<status>}. Kind is the policy
% classification (clarification, local-correction, scope-restriction,
% unfaithful, false-source, open-gap); status is open, wip (elimination
% underway), resolved, or historical. Machine-read by the site index;
% prints a verdict line.
\newcommand{\gapnote}[2]{\par\noindent\textbf{Verdict:} #1 (#2).\par}


\title{Conditional After-Blocking Fundamental Theorem versus the CPSV Statement}
\author{}
\date{2026-05-01}

\begin{document}
\maketitle
\gapnote{open-gap}{open}

\section{The source statement}

The non-periodic matrix product vector Fundamental Theorem in
Refs.~\cite{Cirac2016MPDO_arXiv,Cirac2017MPDO_AnnPhys,Cirac2021Matrix} does not
assume the auxiliary data required by the conditional \Lean{} theorem. The
sources state the comparison only after the tensors have been placed in
canonical form and bases of normal tensors have been selected. This note
records the distinction for the after-blocking theorem in
\href{https://github.com/LionSR/TNLean/pull/1104}{PR \#1104}, discussed in
\href{https://github.com/LionSR/TNLean/issues/1093}{issue \#1093}.

Let $A=(A^i)_{i=1}^{d}$ be a tensor of bond dimension $D$. The 2016 source first
removes invariant subspaces and replaces $A$ by block-diagonal matrices of the
form
\begin{align}
    A^i
        &=
        \bigoplus_{k=1}^{r}\mu_k A_k^i.
    \label{eq:rmp_conditional_after_blocking_ft_statement_canonical_blocks}
\end{align}
These matrices generate the same positive-length matrix product vector family
as the original tensor. Zero blocks are permitted because the retained block
dimensions $D_k$ need satisfy only
\begin{align}
    \sum_{k=1}^{r}D_k
        &\leq D.
    \label{eq:rmp_conditional_after_blocking_ft_statement_retained_dimension}
\end{align}
The construction appears in
Ref.~\cite[lines~214--219]{Cirac2016MPDO_arXiv}. Blocking by a common multiple
of the peripheral periods then removes the periodic components
\cite[lines~227--231]{Cirac2016MPDO_arXiv}. The source concludes that every
tensor has, after blocking, a canonical-form representative with the same
matrix product vector family
\cite[lines~249--251]{Cirac2016MPDO_arXiv}.

The review presents the same two-step argument: it first constructs canonical
form without changing the matrix product vector family and then states the
Fundamental Theorem for tensors already in canonical form
\cite[lines~1752--1758]{Cirac2021Matrix}. It also describes the direct-sum
canonical form and the period-removing blocking
\cite[lines~1798--1820]{Cirac2021Matrix}.

The source comparison uses bases of normal tensors. Such a basis is a minimal
family of normal tensors whose matrix product vectors span those of the original
tensor and are eventually linearly independent
\cite[lines~271--280]{Cirac2016MPDO_arXiv}. If $A_j$ is a basis tensor and
$\mu_{j,1},\ldots,\mu_{j,r_j}$ are the weights of the blocks represented by
$A_j$, the resulting expansion is
\begin{align}
    \ket{V^{(N)}(A)}
        &=
        \sum_j
        \left(\sum_{q=1}^{r_j}\mu_{j,q}^{N}\right)
        \ket{V^{(N)}(A_j)}.
    \label{eq:rmp_conditional_after_blocking_ft_statement_bnt_expansion}
\end{align}
The 2016 paper gives this basis construction and expansion in
Ref.~\cite[lines~271--302]{Cirac2016MPDO_arXiv}; the review gives the same
notion and formula in Ref.~\cite[lines~1846--1885]{Cirac2021Matrix}.

Injectivity is obtained separately. The 2016 source defines injective and
block-injective canonical forms and invokes the quantum Wielandt theorem to
show that normal tensors become injective after additional blocking
\cite[lines~317--344]{Cirac2016MPDO_arXiv}. The review states that blocking at
most $3D^5$ sites gives block-injective canonical form
\cite[lines~2114--2124]{Cirac2021Matrix}.

The source Fundamental Theorem is concise. For canonical-form tensors $A$ and
$B$, with bases of normal tensors $(A_{k_a})_{k_a}$ and $(B_{k_b})_{k_b}$,
proportionality of their matrix product vector families implies that the two
bases have the same cardinality and match, after a permutation, up to phases
and invertible gauge matrices
\cite[Theorem~\texttt{thm1}, lines~347--352]{Cirac2016MPDO_arXiv}. In the equal
case, the source concludes that there is a global gauge
\cite[Corollary~\texttt{II\_cor2}, lines~354--360]{Cirac2016MPDO_arXiv}. The
review repeats both statements
\cite[Theorem~\texttt{thm1} and Corollary~\texttt{II\_cor2},
lines~1891--1900]{Cirac2021Matrix}. It also notes that a later periodic
Fundamental Theorem avoids period-removing blocking
\cite[lines~1908--1909]{Cirac2021Matrix}; the formal derivation considered here
does not use that theorem.

\section{The conditional \Lean{} theorem}

The theorem added in
\href{https://github.com/LionSR/TNLean/pull/1104}{PR \#1104} begins with tensors
$A$ and $B$ having the same formal matrix product vector family. After a
positive blocking, it produces sector decompositions $P$ and $Q$ equipped with
basis-of-normal-tensor data. Its conclusion has the form of the
sector-comparison part of the CPSV statement: the blocked tensors agree with
$P$ and $Q$ at positive lengths, $P$ and $Q$ have the same full matrix product
vector family, and their sector multiplicities and weights match up to a
permutation and nonzero phases.

The formal development no longer uses the former blocked-normal-form records
that packaged this boundary. The remaining auxiliary hypothesis is the
comparison field for the common primitive nonzero-sector families produced by
the preceding structural theorem. The formal theorem derives the blocked-word
relabeling assertion for cyclic-sector data.

The comparison field receives trace-preserving, primitive, and irreducible
structural data, together with positive bond dimensions, for the produced
sector families. It returns basis-of-normal-tensor cover data containing the
zero-tail information, injectivity, normal-form separation, and
block-permutation gauge-phase comparison required for those same families.

Neither Theorem~\texttt{thm1} nor Corollary~\texttt{II\_cor2} states these data
as independent hypotheses
\cite[lines~347--360]{Cirac2016MPDO_arXiv}. In the formal proof they mark the
boundary between the canonical-form construction already established in
\Lean{} and the source's compressed passage from equality of total matrix
product vectors to basis-of-normal-tensor comparison.

\section{Paper-level inputs still to be formalized}

Two components of the comparison field are substantive paper-level inputs. The
sources treat them as part of the theorem or as standard cited input, but the
current \Lean{} proof has not yet derived them for the common-sector families
without invoking the periodic Fundamental Theorem.

First, the injectivity assumptions on the produced nonzero-sector blocks
correspond to the quantum-Wielandt and block-injectivity step in the source
\cite[lines~317--344]{Cirac2016MPDO_arXiv}. The source does not assume
injectivity in Theorem~\texttt{thm1}; it obtains injective or block-injective
tensors after further blocking
\cite[lines~347--352]{Cirac2016MPDO_arXiv}. The formal proof must show that this
refinement applies to the final common primitive sector families used by the
overlap-rigidity comparison.

Second, the basis-of-normal-tensor cover contains the source comparison itself,
specialized to the produced nonzero-sector families. Once canonical forms and
bases of normal tensors have been chosen, the source derives a matching up to
phases and gauges from proportionality or equality of all matrix product
vectors
\cite[lines~347--360]{Cirac2016MPDO_arXiv}. The theorem in
\href{https://github.com/LionSR/TNLean/pull/1104}{PR \#1104} still assumes the
corresponding data for the specific common-sector families: normal
canonical-form structures on both sides, separation of distinct sectors up to
gauge-phase equivalence, zero-tail and injectivity data, and block-matching
data.

A proof that avoids the periodic Fundamental Theorem must derive all these
properties from equality of the original matrix product vector families after
every blocking and reindexing has been transported. They are not additional
assumptions proposed by the papers; they are source-level steps not yet
connected to the present formal construction.

\section{Coordinate and transport conditions}

The remaining hypotheses expose coordinate and transport data that the source
notation identifies implicitly.

The blocked-word relabeling theorem compares two constructions of the common
blocked tensor. One blocks once at the common length; the other first removes
each peripheral period and then groups the resulting blocks. In paper notation,
both constructions are tensors over words of the same blocked length. The
formal theorem derives their coordinate identification from the grouped-block
cast construction.

The former period-removal bridge separately exposed a positive-definite fixed
point of the adjoint transfer map, the primitive root that parametrizes the
peripheral eigenvalues, and irreducibility of the associated completely
positive map. These analytic data are now derived from trace preservation and
tensor irreducibility by the quantum Perron--Frobenius theorem in
\leanid{MPSTensor.exists\_cyclic\_sector\_decomp\_after\_blocking\_of\_TP\_of\_isIrreducibleTensor}.
The blueprint period-removal theorem can therefore state the cyclic-sector
decomposition in a source-faithful form. The remaining gaps concern the later
injectivity refinement, basis-of-normal-tensor cover comparison, zero-tail
bookkeeping, and coordinate transports for the final common-sector families.

The zero-tail condition arises from the distinction between positive-length and
empty-word conventions. The source defines matrix product vectors only at
positive lengths. The formal relation \leanid{SameMPV}\textsubscript{2} also
includes the empty word. A zero block with bond dimension $D_0$ contributes
zero at every positive length but contributes $D_0$ at length zero. The formal
proof must therefore record the zero-tail dimension and either identify the two
zero tails or keep the length-zero equality separate. This condition reflects
the formal convention, not a hypothesis in the source theorem.

Weights and zero-tail identities must also be transported through the selected
common-sector representatives. The comparison theorem applies to the final
block family obtained after common blocking and the canonical relabeling from
the structural construction. The papers move among these descriptions without
naming each transport map; the formal proof must show that the weights,
nonzero-sector tensors, and zero-tail identities are preserved.

\section{Conclusion}

The papers do not state the blocked-normal-form comparison boundary as a list of
theorem hypotheses. The theorem in
\href{https://github.com/LionSR/TNLean/pull/1104}{PR \#1104} is therefore not
yet the unconditional CPSV/CPGSV theorem. It is the result currently proved
without the periodic Fundamental Theorem, with the remaining paper-level inputs
and formal coordinate transports exposed.

The gap lies in the present \Lean{} development, not in the cited sources. The
papers compress the canonical-form reduction, blocking, basis-of-normal-tensor
matching, and injectivity refinement in standard paper notation. The formal
theorem must remain conditional until it proves, for the same final
common-sector data, the injectivity refinement, the basis-of-normal-tensor
cover comparison, and the zero-tail and weight transports. The analytic
period-removal input is now derived and is no longer an open
source-faithfulness gap.

\bibliographystyle{alpha}
\bibliography{references}

\end{document}
